% =============================================================================
% 6.3 研究前景与未来工作
% 草稿版本：v0.2 (2026-05-07)
%   v0.1：作为原 5.4.2 \subsubsection*{未来工作} 起草。
%   v0.2：按导师意见 3 拆出，提级为 6.3 \subsection。\label{sec:future-work}
%        从原 5.4 \subsection 迁到本节，所有外部 \ref{sec:future-work}
%        引用自动指向新位置。原 5.4 节首段"按局限性与未来工作两条线索
%        展开"的过渡话已迁入 6.2 节末段。
% 字数：约 900 字
% 风格规则：v0.2 风格（无 \textbf / \textit / 引例括号 / 双引号 / 破折号）
% 引用：\cite{ref1, ref19, ref29}
% 图表：无
% 依赖宏包：booktabs / tabularx（主稿已加载）
% =============================================================================

\subsection{研究前景与未来工作}
\label{sec:future-work}

针对 \ref{sec:limitations} 节归纳的四类局限性，本节按一一对应的关系
规划 4 类具体的未来工作方向，并在末段给出超出本文研究范围的多智能体
协作扩展作为开放课题。

第一类是显式对话状态跟踪机制，针对多轮长程依赖能力不足。在
\texttt{recognizer} 与 \texttt{completer} 之外引入独立的
\texttt{DialogState} 数据结构，每轮对话结束后把当前 \texttt{locations}、
\texttt{time}、\texttt{intent} 三个槽位写入状态对象，下一轮 prompt 中
以结构化 JSON 形式注入。要求大语言模型对每个槽位显式输出继承或覆盖
的二选一决策，把当前的隐式推理转化为显式标注。同时把多轮评测集从
目前的 3 条扩展到 30 条以上，覆盖意图切换、地点切换、时间切换、复合
切换等典型对话状态切换案例 \cite{ref29}。

第二类是知识库扩展到行业气象领域，针对数值阈值精确归属对分级器规则
的依赖。本文当前知识库覆盖 GB/T 28592、GB/T 28591、GB/T 27964、
GB 50736、GB/T 21987 等 17 项标准共 62 条条目，主要面向公众气象服务
场景。未来将扩展到航空气象 ICAO Annex 3、农业气象作物物候期分级、
海洋气象浪高等级等行业标准，对应同步扩展通路 B 分级器规则与通路 A
评测集。同时引入更细粒度的元数据，例如 \texttt{applicable\_\bk{}region}
区域适用性、\texttt{publish\_\bk{}year} 发布年份等，支持随时间漂移的
标准更新。

第三类是受限沙箱升级为进程级隔离，针对沙箱安全性局限。具体方案有
两条候选路径。一是 RestrictedPython 加 PyPy 沙箱，可以在不引入容器
开销的前提下实现更严格的字节码级访问控制。二是 Pyodide WebAssembly
运行时，把代码执行隔离到 WASM 沙箱中，天然不暴露任何系统调用接口。
两条路径在延迟、安全性、可移植性三个维度上各有取舍，需要进一步实测
对比。

第四类是引用强制注入与生成后校验，针对模块层与集成层之间的指标衰减。具体改进包含
两个方向。一是在 RAG 工具返回结果后立即把 \texttt{must\_\bk{}cite} 关键字
以结构化 JSON 形式注入到模型上下文中，并在 system prompt 中强制要求
最终报告包含该关键字。二是在报告生成完成后引入独立的引用校验后处理器，
对照 \texttt{must\_\bk{}cite} 进行字符串级匹配，未命中时触发模型二次生成。
这两类机制的有效性需要在扩展评测集上独立验证。

进一步的扩展方向包括把单智能体扩展为多智能体协作架构 \cite{ref1, ref19}，
将检索专家、桥接专家、决策专家、报告专家解耦为独立 agent，引入
agent 间共享的 \texttt{Blackboard} 状态空间。多智能体方案在分工
专业化与可调试性上具有优势，但需要解决状态同步与额外延迟两类工程
问题。这一方向超出本文研究范围，保留为后续工作的开放课题。
